\centerline{\bf \Huge Теоретико-числовая диагностика}

\begin{flushright}
        \scriptsize{это должен уметь каждый восьмиклассник}
\end{flushright}

\noindent

\begin{enumerate}

\item Натуральные числа $m$ и $n$ удовлетворяют равенству
\[
    \text{НОД}(m,n)+ \text{НОК}(m,n)=m+n.
\]
Докажите, что одно из чисел $m,n$ делится на другое.
% Источник: ВсОШ, окружной этап, 1995, 10 класс, задача 2.
% https://mathus.ru/math/nodinok.pdf

\item Найдите все пары натуральных чисел $(x,y)$, удовлетворяющие
уравнению
\[
    5xy+y-5x=1038.
\]
% Источник: олимпиада «Покори Воробьёвы горы!», 2014,
% 10--11 классы.
% https://mathus.ru/math/ez.pdf

\item Найдите все простые числа $p$, для которых число
\[
    p^2+11
\]
имеет ровно шесть натуральных делителей.
% Источник: ВсОШ, окружной этап, 1995, 9 класс, задача 5.
% https://mathus.ru/math/primes.pdf

\item Известно, что натуральные числа $m$ и $n$ взаимно просты.
Найдите наибольшее возможное значение числа
\[
    \text{НОД}(m+2000n,\;n+2000m).
\]
% Источник: Московская математическая олимпиада, 2000,
% 11 класс, задача 1.
% https://mathus.ru/math/nodinok.pdf

\item Пусть $m$ и $n$ — натуральные числа. Докажите, что
\[
    (2^m-1)^2\mid 2^n-1
\]
тогда и только тогда, когда
\[
    m(2^m-1)\mid n.
\]
% Источник: ВсОШ, окружной этап, 1997, 10 класс, задача 3.
% https://mathus.ru/math/ostatki.pdf

\item Найдите все натуральные числа $n$, для которых существуют
натуральные числа $a$ и $b$, удовлетворяющие равенству
\[
    n!=2^a+2^b.
\]

% Источник: муниципальный этап ВсОШ в Москве,
% 2020/21 учебный год, 11 класс, задача 5.
% https://mmo.mccme.ru/2020/okrug/Resh_11.pdf

\item Три натуральных числа $a,b,c$ удовлетворяют условиям
\[
    a+b\mid ab,\qquad
    b+c\mid bc,\qquad
    c+a\mid ca.
\]
Докажите, что числа $a,b,c$ имеют общий делитель, больший единицы.
% Источник: ВсОШ, окружной этап, 2004,
4% 9 класс, задача 4; 10 класс, задача 3.
% hраttps://mathus.ru/math/nodinok.pdf

\item Натуральные числа $a,b,c$, где $c>2$, удовлетворяют равенству
\[
    \frac1a+\frac1b=\frac1c.
\]
Докажите, что хотя бы одно из чисел
\[
    a+c,\qquad b+c
\]
является составным.
% Источник: региональный этап ВсОШ, 2013,
% 10 класс, задача 6.
% https://mathus.ru/math/nodinok.pdf

\end{enumerate}